% Bab 3: Product Requirements - Baseline
\chapter{Product Requirements - Baseline}
\label{bab:product-requirements-baseline}
\begin{longtable}{|L{0.075}|L{0.226}|L{0.299}|L{0.093}|L{0.307}|}
\caption{Product Requirements - Baseline}\\
\hline
\cellcolor[HTML]{D9EAF7}\textbf{ID} & \cellcolor[HTML]{D9EAF7}\textbf{Requirement} & \cellcolor[HTML]{D9EAF7}\textbf{Deskripsi Implementasi} & \cellcolor[HTML]{D9EAF7}\textbf{Priority} & \cellcolor[HTML]{D9EAF7}\textbf{Acceptance Criteria} \\ \hline
\endfirsthead
\hline
\cellcolor[HTML]{D9EAF7}\textbf{ID} & \cellcolor[HTML]{D9EAF7}\textbf{Requirement} & \cellcolor[HTML]{D9EAF7}\textbf{Deskripsi Implementasi} & \cellcolor[HTML]{D9EAF7}\textbf{Priority} & \cellcolor[HTML]{D9EAF7}\textbf{Acceptance Criteria} \\ \hline
\endhead
F01 & Location Awareness & Aplikasi membaca posisi pengguna (akurasi tinggi, pembaruan tiap \(\pm\)3 detik atau perpindahan \(\ge\)5 m). Posisi ditampilkan sebagai titik biru dengan lingkaran akurasi di peta, disertai status ``Di dalam / di luar area Sudirman Street Food''. Jika izin lokasi ditolak, aplikasi berjalan dalam mode tanpa lokasi: peta dan daftar tenant tetap dapat dibuka secara manual, tetapi trigger otomatis dan navigasi tidak aktif. & Must & (1) Posisi tampil di peta \(\le\)10 detik setelah aplikasi dibuka dengan GPS aktif.\newline (2) Posisi diperbarui saat pengguna berjalan \(\ge\)5 m.\newline (3) Status di dalam/di luar area berubah dengan benar saat pengguna melewati batas area.\newline (4) Jika izin ditolak, muncul layar penjelasan dan konten tenant tetap dapat dibuka manual dari peta atau daftar. \\ \hline
F02 & Location-Based Trigger (Geofence) & Trigger sepenuhnya berbasis geofence GPS dengan empat jenis event:\begin{itemize}\item Entry event: geofence kawasan radius 100 m memunculkan notifikasi sambutan.\item Proximity event: geofence tenant radius 10 m (nilai \texttt{triggerRadiusM} per POI pada dataset) memunculkan kartu POI dan tawaran audio. Dicek dari pembaruan lokasi background dengan perhitungan jarak (haversine).\item Dwell event: pengguna berada di dalam geofence tenant \(\ge\)3 menit, sehingga kunjungan tercatat otomatis; jika melebihi alokasi trail, muncul pengingat sisa waktu (F08).\item Exit event: pengguna keluar dari geofence tenant, lalu muncul rekomendasi berikutnya (F07).\end{itemize}Pencegahan trigger palsu/berulang: pembacaan lokasi dengan akurasi \textgreater{}30 m diabaikan, hysteresis (exit tercatat jika jarak \textgreater{} radius + 10 m), dan cooldown 30 menit per POI. Trigger berjalan di background melalui foreground service selama ``Mode Jelajah'' aktif. & Must & (1) Notifikasi entry muncul \(\le\)15 detik setelah pengguna masuk area, termasuk saat aplikasi di background atau layar terkunci.\newline (2) Kartu POI muncul \(\le\)10 detik saat pengguna masuk radius POI, berhasil minimal 8 dari 10 percobaan lapangan.\newline (3) Trigger tidak muncul berulang selama pengguna tetap di area POI.\newline (4) Kunjungan tercatat setelah pengguna berada \(\ge\)3 menit di geofence tenant; pengguna yang hanya lewat (\textless{}3 menit) tidak tercatat berkunjung.\newline (5) Pembacaan lokasi dengan akurasi \textgreater{}30 m tidak memicu trigger. \\ \hline
F03 & Multimodal Exploration & Dataset memuat 31 POI tenant di Sudirman Street Food. Setiap POI memiliki: foto tenant/makanan (jika tersedia), kategori dan profil rasa, daftar menu dan harga, jam buka, serta audio narasi 45--90 detik (TTS) yang dilengkapi transkrip. Cerita kuliner/budaya (asal-usul makanan, keunikan tenant) disampaikan melalui narasi audio dan transkripnya. Seluruh konten (data JSON, foto, dan audio) dibundel di dalam APK sehingga dapat diakses offline. & Must & (1) 31/31 POI memiliki menu/harga, jam buka, serta audio narasi beserta transkrip.\newline (2) Audio dapat diputar, dijeda, dan dilanjutkan tanpa koneksi internet.\newline (3) Transkrip audio tersedia untuk setiap narasi.\newline (4) Konten POI terbuka \(\le\)2 detik. \\ \hline
F04 & Eyes-Free Interaction & Saat trigger terjadi, aplikasi memutar prompt suara, misalnya ``Kamu mendekati Batagor X. Mau dengar ceritanya?'', disertai getar. Pengguna merespons lewat tombol aksi di notifikasi (Dengar / Lewati) tanpa membuka aplikasi. Mode Jalan memutar narasi secara otomatis. Navigasi memberi panduan audio (``Tenant tujuan 20 meter lagi di sebelah kanan'') dan pengumuman saat tiba di tujuan. & Must & (1) Pengguna dapat memutar narasi dari notifikasi lock screen tanpa membuka aplikasi.\newline (2) Dalam Mode Jalan, narasi mulai \(\le\)3 detik setelah trigger.\newline (3) Panduan audio navigasi berbunyi minimal setiap pengurangan jarak 20 m dan saat tiba.\newline (4) Dalam usability test, minimal 4 dari 5 penguji dapat menyelesaikan satu kunjungan POI (navigasi, trigger, dan narasi) tanpa melihat layar. \\ \hline
F05 & Search \& Destination & Pencarian berdasarkan nama tenant, nama menu, atau kategori, dengan pencarian toleran typo. Filter yang tersedia: rentang harga (\textless{}Rp15rb, Rp15--30rb, \textgreater{}Rp30rb), rasa (pedas/manis/gurih), jenis (makanan/minuman/jajanan), dan status buka. Hasil diurutkan berdasarkan jarak, dan pengguna dapat menetapkan tenant sebagai tujuan. & Must & (1) Hasil pencarian muncul \(\le\)1 detik secara offline.\newline (2) Kata kunci ``batagor'' dan typo ``batagr'' sama-sama menampilkan tenant batagor.\newline (3) Kombinasi filter menghasilkan daftar yang sesuai.\newline (4) Tenant dapat dipilih sebagai tujuan dan langsung masuk ke mode navigasi. \\ \hline
F06 & Navigation & Menuju kawasan: membuka Google Maps melalui intent dengan koordinat kawasan.\newline Di dalam kawasan: peta (MapLibre dengan tile OpenStreetMap) menampilkan rute jalan kaki dari posisi pengguna ke tenant tujuan. Rute dihitung oleh layanan OSRM publik (profil pejalan kaki) dan dihitung ulang jika pengguna keluar \textgreater{}25 m dari rute; jika layanan gagal atau tidak ada koneksi, aplikasi memakai garis lurus. Peta juga menampilkan jarak dalam meter, dan panah arah berbasis kompas. Kedatangan terdeteksi saat jarak \(\le\)10 m, lalu aplikasi memutar pengumuman ``Kamu sudah sampai!'' dan membuka kartu POI tenant tujuan. & Must & (1) Jarak ke tujuan diperbarui saat pengguna berjalan.\newline (2) Panah menunjuk arah tujuan dengan deviasi \(\le\)30\textdegree{}.\newline (3) Notifikasi kedatangan muncul saat pengguna berada \(\le\)10 m dari tenant, dan kartu POI tenant tersebut terbuka.\newline (4) Navigasi dapat dilakukan ke 31/31 POI.\newline (5) Tombol ``Buka di Google Maps'' berfungsi dari luar kawasan.\newline (6) Tanpa koneksi internet, navigasi tetap berjalan dengan garis lurus ke tujuan. \\ \hline
F07 & Contextual Recommendation & Rekomendasi berbasis aturan (rule-based scoring) dengan faktor: jarak, status buka, kecocokan preferensi rasa, sisa budget, jam ramai (data observasi), dan status sudah/belum dikunjungi (dari dwell event). Muncul saat pengguna keluar dari geofence tenant yang sudah dikunjungi (exit event) dalam bentuk ``3 pilihan berikutnya'', masing-masing disertai alasan (misalnya ``30 m dari sini \(\cdot\) sesuai budget \(\cdot\) biasanya tidak ramai jam segini''). Ada juga aturan pasangan rasa, misalnya setelah makanan pedas disarankan minuman segar. & Must & (1) Rekomendasi tidak pernah memuat tenant yang tutup atau sudah dikunjungi.\newline (2) Setiap rekomendasi menampilkan alasannya.\newline (3) Rekomendasi muncul \(\le\)1 detik setelah exit event.\newline (4) Input yang sama selalu menghasilkan rekomendasi yang sama (deterministik). \\ \hline
F08 & Domain-Specific Innovative Feature: Smart Culinary Trail Planner & Pengguna mengisi waktu (30/60/90 menit), budget, dan preferensi rasa. Aplikasi memilih kombinasi tenant yang sesuai, lalu menyusun urutan rute dengan algoritma nearest neighbor, lengkap dengan estimasi total biaya dan waktu (dari data harga dan rata-rata durasi kunjungan hasil observasi). Progres trail diperbarui otomatis dari dwell event. Selama trail berjalan, dwell event juga memberi pengingat jika pengguna melewati alokasi waktu di satu tenant, misalnya ``Sisa waktumu 25 menit, masih ada 2 tenant.'' & Must & (1) Untuk input 60 menit dan Rp50.000, aplikasi menghasilkan itinerary minimal 2 tenant dengan estimasi biaya \(\le\) budget dan waktu \(\le\) durasi.\newline (2) Urutan rute tidak bolak-balik (setiap langkah menuju tenant terdekat berikutnya).\newline (3) Tenant pada trail otomatis ditandai selesai setelah dwell event tercatat.\newline (4) Notifikasi sisa waktu muncul saat dwell melebihi alokasi + 5 menit.\newline (5) Pengguna dapat melewati tenant atau menyusun ulang trail. \\ \hline
F09 & Location Privacy \& Consent & Layar onboarding menjelaskan alasan aplikasi memakai lokasi, kapan lokasi dibaca (hanya saat Mode Jelajah aktif), dan bahwa riwayat lokasi tidak dikirim ke server aplikasi. Riwayat posisi, kunjungan, dan progres trail hanya disimpan di perangkat. Satu-satunya pengiriman koordinat adalah saat navigasi aktif (F06): titik asal dan tujuan dikirim ke layanan rute OSRM publik untuk menghitung rute jalan kaki, tanpa identitas pengguna. Pengguna dapat mematikan Mode Jelajah kapan saja. & Must & (1) Penjelasan tampil sebelum permintaan izin sistem dan menyebut pemakaian layanan rute saat navigasi.\newline (2) Pelacakan background berhenti saat Mode Jelajah dimatikan (notifikasi foreground hilang).\newline (3) Tidak ada request jaringan ke server aplikasi yang membawa data lokasi; request yang membawa koordinat hanya ke layanan rute OSRM saat navigasi aktif (diverifikasi lewat log). \\ \hline
\end{longtable}

\section{Additional Requirements (Optional)}

\begin{longtable}{|L{0.082}|L{0.224}|L{0.296}|L{0.093}|L{0.304}|}
\caption{Additional Requirements (Optional)}\\
\hline
\cellcolor[HTML]{D9EAF7}\textbf{ID} & \cellcolor[HTML]{D9EAF7}\textbf{Requirement} & \cellcolor[HTML]{D9EAF7}\textbf{Deskripsi} & \cellcolor[HTML]{D9EAF7}\textbf{Priority} & \cellcolor[HTML]{D9EAF7}\textbf{Acceptance Criteria} \\ \hline
\endfirsthead
\hline
\cellcolor[HTML]{D9EAF7}\textbf{ID} & \cellcolor[HTML]{D9EAF7}\textbf{Requirement} & \cellcolor[HTML]{D9EAF7}\textbf{Deskripsi} & \cellcolor[HTML]{D9EAF7}\textbf{Priority} & \cellcolor[HTML]{D9EAF7}\textbf{Acceptance Criteria} \\ \hline
\endhead
F10 & Paspor Kuliner (Gamification) & Setiap kunjungan yang tercatat lewat dwell event (\(\ge\)3 menit di geofence tenant) memberi stempel digital; tidak ada scan QR. Tersedia trail tematik, misalnya ``Jejak Jajanan Sunda'' dan ``Manis-manis Bandung''. Badge diberikan saat satu trail selesai. Data disimpan lokal (AsyncStorage). & Could & (1) Stempel tercatat setelah dwell event dan tidak terduplikasi pada tenant yang sama.\newline (2) Progres trail tampil (misalnya 3/5).\newline (3) Badge muncul saat trail selesai.\newline (4) Data tetap ada setelah aplikasi ditutup dan dibuka kembali. \\ \hline
F11 & Akun Pengguna Lokal & Pengguna masuk atau mendaftar (email, nama pengguna, kata sandi) sebelum onboarding. Sesi disimpan di perangkat melalui SecureStore dan tidak dikirim ke server. Tersedia opsi ``Ingat Saya'' dan halaman Profil untuk melihat data akun, mengubah preferensi, dan keluar. & Could & (1) Pengguna yang belum masuk diarahkan ke halaman login.\newline (2) Setelah masuk, pengguna lanjut ke onboarding (pertama kali) atau Beranda.\newline (3) Sesi tetap tersimpan setelah aplikasi dibuka ulang.\newline (4) Keluar akun menghapus sesi dari perangkat. \\ \hline
\end{longtable}
